\chapter{\abstractname}

%TODO: Abstract
Music plays a key role in everyday life, entertaining and engaging complex cognitive and emotional processes in the human brain. Nonetheless, the neural mechanisms behind music perception remain unknown. Therefore, this thesis explores whether \ac{EEG} signals can be used to reliably decode music perception, aiming to bridge the gap between \ac{MIR} and neurosciences. A dedicated dataset for \ac{EEG} based Music Genre classification is created to contribute to this area, named NeuroMGR. By utilising the GTZAN music dataset, \ac{EEG} signals have been recorded by a hybrid biosignal headset while listening to each music clip. The recorded signal then underwent a custom denoising process to eliminate most artifacts, followed by an exhaustive feature extraction combined with specialised feature selection methods. Based on these features, machine learning classifiers were trained to predict the corresponding features, which are then analysed regarding accuracy, F1-score, and their confusion matrices. To compare with a state-of-the-art \ac{EEG} dataset, the same approach is followed with the \ac{DEAP} dataset, which is evaluated in the same way as the custom dataset. Results demonstrate that genre-specific neural patterns can be detected above chance level. Concreteley, mean accuracies of 68.3\% and 20.2\% were measured for binary classification on the \ac{DEAP} dataset and ten-class classification on the NeuroMGR dataset respectively. Although signal quality and device limitations posed a challenge, the key findings are promising evidence that \ac{EEG} signals possess meaningful information about musical perception. Overall, this thesis lays the groundwork for further research in the research area of \ac{EEG}-based \ac{MIR} by providing a novel dataset alongside detailed methodology for \ac{MGR} using \ac{EEG} signals. Future studies can build upon this thesis by employing larger subject pools, advanced deep learning architectures, or multimodal data fusion. 
